\section{Evaluating Both Boundaries}\label{sec:reporting}
An intervention needs one record linking the attacker's remaining observations to the original task and allowed collateral effects. Table~\ref{tab:availability} shows why this joint record matters: among four query/workflow privacy settings from two sources, adaptation is documented in two and matched original-task utility in one, with no setting documenting both. The adapted settings share a preprint; the utility observation is WRA's published visibility comparison. These are descriptive counts of inspected configurations.

\begin{table*}[t]\centering\small
\caption{Evidence located in the retained records. Counts use eligible settings; the second column also gives distinct sources. ``Not located'' leaves evidence unresolved within the inspected scope. These counts describe availability, not study quality or checklist validity.}\label{tab:availability}
\setlength{\tabcolsep}{4pt}
\begin{tabularx}{\textwidth}{@{}P{.28\textwidth}P{.15\textwidth}P{.11\textwidth}Y@{}}\toprule
Question & Eligible settings / sources & Located & Supporting settings and consequence \\\midrule
Protected-condition attack adaptation & 4 / 2 & 2/4 & Communication graphs: transformed captures and live OHTTP. Both come from one preprint. \\\addlinespace[3pt]
Original-task utility under the same intervention & 4 / 2 & 1/4 & WRA visibility comparison. It does not price WRA's executed-decoy condition. \\\addlinespace[3pt]
Adaptation and original-task utility in one setting & 4 / 2 & 0/4 & Adaptation and utility are located in different records; transport cost alone does not fill the task-outcome field. \\\addlinespace[3pt]
Original-task utility for executed cover & 1 / 1 & 0/1 & WRA's five research decoys report recovery and extra research; matched original-task utility is not located. \\\bottomrule
\end{tabularx}
\end{table*}

\begin{table*}[t]\centering\small
\caption{Testable open questions derived from the synthesis. Q1 specifies transfer tests. Q2 separates the forwarding mediators; Q3 extends the controlled cover study to application constraints; Q4 extends live detection to new workloads and deployments. Each gives a result that could support or revise the proposed explanation.}\label{tab:open-questions}
\setlength{\tabcolsep}{4pt}
\begin{tabularx}{\textwidth}{@{}P{.22\textwidth}P{.43\textwidth}Y@{}}\toprule
Question and unresolved mechanism & Controlled comparison & Discriminating result \\\midrule
\textbf{Q1. Which signals survive execution changes?} Do attacks identify recurring task structure or features tied to its implementation? & Keep task labels and observation rights fixed. Hold out compositions, coordination and client/scaffold versions separately, then jointly. Compare destination, volume/timing, sequence and recipient-visible feature families under declared frozen and adapted training. & Identify a feature family whose inference survives specified changes, or a failure tied to one changed condition. Report per-axis and joint transfer with attempted-run coverage; a pooled average cannot establish invariance. \\\addlinespace[4pt]
\textbf{Q2. How does feedback reach model traffic?} Which part is content, length or later action structure? & Fix initial instructions; project to category-free task fields. Pair length-matched neutral padding with validity-preserving replacement of supplementary content; hold encoding and follow-up policy fixed where possible. Vary record count separately. Compare prefix, first-feedback and complete model views under tool-only and joint protection. & Additional held-out inference after feedback, with controls that exclude the competing path, supports that attribution. Report an effect interval against a prespecified meaningful-effect threshold. A nonsignificant difference alone does not rule out leakage. \\\addlinespace[4pt]
\textbf{Q3. Can cover preserve task semantics?} When should it stop at transport or execute in isolation? & Compare planner-sampled decoy schedules (destination sequence and timing) with marginally matched dispatch under production tool contracts and legitimate concurrent updates. Compare no cover, transport termination and isolated execution under declared resource budgets. Include timing-informed removal and independently captured record times; measure inference, delay, original-task completion, quota, state and planner-visible feedback. & A privacy--utility operating point satisfying prespecified state and task tolerances supports the design. Resource exhaustion, changed business state or impaired completion identifies the constraint it violates. \\\addlinespace[4pt]
\textbf{Q4. Which inference survives task discovery?} Does cross-link association recover information hidden by waits and multiplexing? & Train on one concurrency/load and test new compositions, loads and gateways. Compare tool-only and associated tool/model views with observable flow identifiers; keep ground-truth task and stream mappings outside attacker inputs. Include real unrelated traffic and constant-output detector references. & Measure interval recall, false alarms per unit time and correctly localized inference together. Gains must survive changed workloads and loads at a stated false-alarm budget. \\
\bottomrule\end{tabularx}\end{table*}


\subsection{A minimal evaluation procedure}
The following procedure connects established evaluation guidance~\cite{arp2022dos,juarez2014critical,mathews2023critical} to the information and action boundaries. It is a proposed reporting procedure, illustrated by the controlled cover case. The six observation-contract fields in Section~3 specify an inference setting; the eight checklist items organize the broader study report, including its outcomes, uncertainty and provenance.

\begin{enumerate}
\item \textbf{Declare the target and permitted observations.} Name the private fact, visible model/tool/recipient interfaces, access rights and any supplied execution windows. Track where tool results are forwarded.
\item \textbf{Declare the allowed execution effects.} Specify the original goal, other objects that may change, resource budgets and what cover may return to the planner. Use business invariants that permit legitimate concurrent updates; exact state equality fits only a correspondingly constrained fixture.
\item \textbf{Pair protection, attacks and outcomes.} Evaluate the baseline and intervention on matched tasks. Train permitted attackers on protected executions and include detection--removal and joint-decoding attacks that exploit cover timing, position or schedule. Record goal completion, completed tasks violating the state contract, and bytes, delay and resource use on those same runs.
\item \textbf{Make the decision against stated tolerances.} Report attack-specific inference and uncertainty together with task, state and resource tolerances fixed before evaluation. A completed task with a state change outside the declared contract fails its utility requirement. Successful inference on another permitted view defeats a claim of containment across those views. An unsuccessful tested attack supplies no general secrecy guarantee.
\end{enumerate}

\paragraph{Worked application}
Section~\ref{sec:cover-study} follows steps 1--3 for a tool-metadata observer. Step 4 is where completion-only reporting would have failed: live writes pass the goal predicate in 102/120 Qwen runs but violate the state contract in 95 of them. Supplement S8 derives the fixture contract; Supplements S9--S11 add the timing-informed residual views.

\subsection{Reusable forms}
Supplement S5 provides the eight-field checklist, its correspondence to established guidance, and completed WRA, OHTTP and ANT forms. Each form records evidence, a missing item or inapplicability. Identity-obfuscation (3), adversarial-Agent blocking (2) and non-protection (35) settings keep separate denominators. The access map assigns 22 transit, 10 recipient, 11 active and one downstream preference setting.
